\documentclass[10pt,a4paper]{amsart}
\usepackage[margin=19mm]{geometry}
\usepackage{amsmath,amssymb,mathtools}
\usepackage[scr=rsfs,cal=euler]{mathalfa}
\usepackage{xfrac}
\usepackage[unicode,hidelinks,bookmarks=false]{hyperref}
\newcommand{\ie}{i.e.\ }
\newcommand{\cf}{cf.\ }
\newcommand{\resp}{respectively\ }
\newcommand{\Subsection}{\subsection}
\providecommand{\nobreakdash}{\nobreak\hbox{-}}
\setlength{\emergencystretch}{2em}
\multlinegap0pt
\usepackage{amssymb}
\newtheorem{proposition}{Proposition}
\newcommand{\F}{\mathbb{F}}
\newcommand{\cY}{\mathcal{Y}}
\title{Maximal and minimal curves of the form $y^3=x^{(q^2+1)/2}+x$}
\author{Guilherme Dias, Saeed Tafazolian}
\date{Source: arXiv:2608.20654v1 (CC BY 4.0)}
\begin{document}
\maketitle

\noindent\emph{Source and licence.} Maximal and minimal curves of the form $y^3=x^{(q^2+1)/2}+x$. Version arXiv:2608.20654v1 (CC BY 4.0), distributed by arXiv under the Creative Commons Attribution 4.0 International licence, http://creativecommons.org/licenses/by/4.0/, as recorded in the arXiv licence metadata for this exact version and shown on the abstract page https://arxiv.org/abs/2608.20654v1. Full licence notice: \texttt{licenses/arxiv-2608.20654-CC-BY-4.0.txt}.\par\smallskip

\noindent\textit{Editorial scope.} Counted unit: the article's proposition prop:count (source0.tex lines 288--295). The article's complete original proof of it is delivered with it (source0.tex lines 297--349, one \texttt{proof} environment of 53 lines). No mathematics is written, repaired or re-derived: the statement and the proof are the article's own lines, in the article's own notation, and no retained line is substituted or re-flowed.  Retained uncounted as the source-local context the proof needs: lines 124--126.  Nothing was omitted from any retained span, so the package carries no omission record.  No retained line contains a citation, so the wrapper carries no bibliography and no bibliography entry.  No retained line contains a figure environment, an image-inclusion command or drawing code, so no image is delivered and the package declares no asset dependency.  Editorial formatting only, in addition: an AMS \texttt{amsart} preamble that restates the article's own declarations and macros used by the retained lines, namely the environments \texttt{proposition} on separate counters, together with the standard packages those lines need.  The retained environments print in this document as Proposition 1 in order of appearance, whereas the article prints the same items with the numbering of its own full document.  The complete native source of the article as distributed in the arXiv e-print for this exact version is bundled unchanged as \texttt{sources/208194/source0.tex}.
\begin{equation}\label{eq:roots}
\#\{x\in\F_Q:x^n=a\}=\sum_{j=0}^{n-1}\chi^j(a),
\end{equation}

\begin{proposition}\label{prop:count}
The smooth projective model of $\cY$ satisfies
\begin{equation}\label{eq:pointcount}
\#\cY(\F_Q)=Q+1+
\sum_{\substack{1\le j\le 3d-1\\3\nmid j}}
J(\chi^j,\chi^{dj}).
\end{equation}
\end{proposition}

\begin{proof}
For the Kummer extension
\[
v^{3d}=t(1-t)^d,
\]
the valuation of the right-hand side at $t=0$ is $1$, so the extension has degree $3d$. The only ramified places of $\F_Q(t)$ are $t=0$, $t=1$, and $t=\infty$, with corresponding valuations
\[
1,\qquad d,\qquad -(d+1).
\]
Thus
\[
\gcd(3d,1)=1,
\qquad
\gcd(3d,d)=d,
\qquad
\gcd(3d,d+1)=1.
\]
Hence there is a unique place above $t=0$ and a unique place above $t=\infty$.

We verify that the $d$ places above $t=1$ are rational. Since $q$ is odd,
\[
q^2\equiv 1\pmod 8,
\]
so $d=(q^2-1)/2$ is even. Put $u=t-1$. Then
\[
t(1-t)^d=u^d(1+u),
\]
and locally at $u=0$,
\[
\left(\frac{v^3}{u}\right)^d=1+u.
\]
After reduction modulo $u$, the residual equation is
\[
Z^d=1.
\]
Since $d\mid Q-1$ and $p\nmid d$, it has exactly $d$ distinct roots in $\F_Q$. Hence there are exactly $d$ $\F_Q$-rational places above $t=1$. Therefore the total contribution above $t=0,1,\infty$ is $d+2$.

For $t\in\F_Q\setminus\{0,1\}$, \eqref{eq:roots} gives
\[
\#\{v:v^{3d}=t(1-t)^d\}
=
\sum_{j=0}^{3d-1}\chi^j(t)\chi^{dj}(1-t).
\]
The term $j=0$ contributes $Q-2$. If $3\mid j$ and $j\ne 0$, then $\chi^{dj}=\varepsilon$ and the inner sum is $-1$; there are $d-1$ such indices. If $3\nmid j$, both $\chi^j$ and $\chi^{dj}$ are nontrivial, and the inner sum is
\[
J(\chi^j,\chi^{dj}).
\]
Indeed, the omission of $t=0,1$ does not alter the Jacobi sum, since our nontrivial characters are extended by zero at the origin. Combining the unramified fibers with the $d+2$ rational places lying above the branch points gives
\[
(Q-2)-(d-1)+(d+2)=Q+1,
\]
which proves \eqref{eq:pointcount}.
\end{proof}

\end{document}
